\chapter{Hydronephrosis}
\index{Hydronephrosis}
\label{sec:Hydronephrosis}

\section{Overview}
Hydronephrosis is the dilation of the renal pelvis and calyces caused by obstruction of urine outflow from the kidney. It can be unilateral or bilateral, acute or chronic, and affects patients of all ages from prenatal to elderly. The condition may be asymptomatic or lead to progressive renal damage if the underlying obstruction is not relieved.
\section{Classification}

\begin{description}[font=\normalfont\bfseries,leftmargin=3em,labelwidth=2.5em]
  \item[Cause] Obstructive / Anatomic
  \item[Organ System] Kidneys
  \item[Pathophysiology] Urinary tract obstruction leading to increased intrapelvic pressure, tubular dilation, and eventual renal parenchymal damage
  \item[Duration] Acute (hours to days) or Chronic (months to years)
  \item[Onset] Gradual or Sudden
  \item[Mode of Transmission] Non-communicable
  \item[Clinical Course] Variable depending on cause and duration
  \item[Severity] Mild to Severe
  \item[Extent of Disease] Localized (unilateral or bilateral)
  \item[Age Group] All ages
  \item[Epidemiology] Common; found in 1--5\% of autopsies; congenital in 1 in 100 births
  \item[ICD-11 Code] GB00.0
\end{description}

\section{Causes}
Hydronephrosis results from any obstruction along the urinary tract, from the ureteropelvic junction to the urethra. Common causes include ureteropelvic junction obstruction (congenital), ureteral stones, benign prostatic hyperplasia, pelvic or retroperitoneal tumors, posterior urethral valves (neonates), and neurogenic bladder. Vesicoureteral reflux can also produce pelvicalyceal dilation without obstruction.

\section{Risk Factors}

\begin{itemize}[noitemsep]
  \item Nephrolithiasis (kidney stones)
  \item Benign prostatic hyperplasia
  \item Pelvic malignancies (prostate, bladder, cervical, ovarian, colorectal)
  \item Pregnancy (physiologic hydronephrosis)
  \item Congenital anomalies (UPJ obstruction, posterior urethral valves)
  \item Neurogenic bladder
  \item Previous pelvic surgery or radiation
  \item Spinal cord injury
\end{itemize}

\section{Signs and Symptoms}

\subsection{Early symptoms}
\begin{itemize}[noitemsep]
  \item Early manifestations of hydronephrosis may be subtle and nonspecific
\end{itemize}

\subsection{Common symptoms}
\begin{itemize}[noitemsep]
  \item \subsection{Early symptoms}
  \item \subsection{Common symptoms}
  \item \textbf{Common symptoms:}
  \item Flank pain or discomfort, often radiating to the groin
  \item Acute obstruction causes severe colicky pain (ureteral colic)
  \item Nausea and vomiting with acute obstruction
  \item Reduced urine output or anuria (bilateral or solitary kidney)
  \item Urinary tract infections and fever
  \item Hematuria (especially with stones)
  \item Hypertension from activation of the renin-angiotensin system
  \item Pain or discomfort in the affected area
  \item Difficulty performing daily activities
\end{itemize}

\subsection{Less common symptoms}
\begin{itemize}[noitemsep]
  \item Atypical or less common presentations of hydronephrosis may occur
\end{itemize}

\subsection{Emergency warning signs}
\begin{itemize}[noitemsep]
  \item Severe or worsening symptoms of hydronephrosis require immediate medical evaluation
\end{itemize}
\section{Progression and Stages}
If obstruction is unrelieved, progressive dilation leads to tubular atrophy, interstitial fibrosis, and irreversible nephron loss. Acute complete obstruction causes rapid renal damage within days. Chronic partial obstruction allows for slow, insidious loss of renal function over months to years. The extent of recovery after relief depends on the severity and duration of obstruction.

\section{Complications}

\begin{itemize}[noitemsep]
  \item Obstructive uropathy and hydronephrotic atrophy
  \item Acute kidney injury or chronic kidney disease
  \item Urinary tract infections and pyelonephritis
  \item Urolithiasis
  \item Hypertension
  \item Sepsis from infected obstructed kidney (pyonephrosis)
  \item Reduced quality of life
  \item Emotional and psychological effects
\end{itemize}

\section{Diagnosis}

\begin{itemize}[noitemsep]
  \item Renal ultrasound is the initial imaging modality of choice
  \item CT urography or MRI to identify level and cause of obstruction
  \item Voiding cystourethrogram (VCUG) for suspected reflux or posterior urethral valves
  \item Renal scintigraphy (MAG3) to assess functional significance
  \item Blood tests (creatinine, BUN) to evaluate renal function
  \item Urinalysis and culture to exclude infection
\end{itemize}

\section{Prevention}

\begin{itemize}[noitemsep]
  \item Maintain hydration to reduce stone formation
  \item Regular urologic screening for patients with known risk factors
  \item Prompt treatment of urinary tract infections
  \item Prenatal ultrasound screening for congenital anomalies
  \item Go for regular check-ups so any problems are caught early
\end{itemize}

\section{Treatment Options}

\begin{itemize}[noitemsep]
  \item Relief of obstruction by ureteral stenting or nephrostomy tube
  \item Ureteroscopy with laser lithotripsy for stones
  \item Pyeloplasty for ureteropelvic junction obstruction
  \item Transurethral resection of prostate (TURP) for BPH
  \item Valve ablation for posterior urethral valves
  \item Antibiotics for associated urinary tract infection
  \item Lifestyle changes and self-care
  \item Surgery when needed
\end{itemize}

\section{Living with It}

\begin{itemize}[noitemsep]
  \item Follow your treatment plan as prescribed
  \item Keep up with regular appointments with your healthcare provider
  \item Adjust your daily habits as needed
  \item Reach out to family, friends, or support groups for help
  \item Keep learning about your condition and how to manage it
\end{itemize}

\section{Outlook}
The prognosis depends on the cause, severity, and duration of obstruction. Acute obstruction relieved promptly carries an excellent prognosis with full recovery of renal function. Chronic or severe obstruction can lead to irreversible renal damage. Long-term monitoring after relief of obstruction is essential to detect recurrence and manage residual renal impairment.
